\scanpage{54}
\mspage{7}
$\eta_0=\operatorname{Coker}(\pi_1\longrightarrow\eta)$, et l'hom
\[
\lambda_0:\pi_0\otimes\pi_0\longrightarrow\eta_0
\]
déduit de $\lambda_G$ par passage au quotient. C'est
indépendant du choix de $G$, on le voit que
c'est ce qu'on construit directement dans
le groupe $K=G/\pi_1$. Il faut donc

b) $\lambda_0$ surjectif (car signifie aussi
que $N/\pi_1$ dans $K\subset M\times M'$ est $=DK$,
ou encore que $dN=DM\ \unsure{\}}$ exactement)
\margin{Vérifier}

\deleted{\unsure{a}} Introduisons alors $\operatorname{Ker}\lambda_0=P\subset\pi_0\otimes\pi_0$,
(indépendant de $G$), donc $\lambda_G$ induit
\[
\overline\lambda_G:P\twoheadrightarrow\pi_1^{\deleted{\unsure{*}}}{}_{\inserted{*}}=\operatorname{Im}(\pi_1\longrightarrow\eta)\subset\eta
\]
et la condition effective dépendant du
choix de $G$ qui reste à exprimer
est

c) $P\xrightarrow{\overline\lambda_G}\pi_1^{\deleted{\unsure{*}}}{}_{\inserted{*}}$ est un épi.

Ici la dépendance de cette condition
par $G$ variable est explicitée grâce à
la formule donnant $\overline\lambda_{\unsure{G_1}}$ en fonction
de $\lambda_G$, qui s'écrit de façon évidente
\[
\overline\lambda_{G_1}=\overline\lambda_G+pc_Ei
\]
$i$ l'inclusion $P\subset\pi_0\otimes\pi_0$, $p$ la projection $\pi_1\longrightarrow\pi_{1*}$.
\mspage{53}
